% Thông tin nhóm, phân công theo mảng việc, công cụ AI; hiện trạng sau tuần đầu chỉ ở báo cáo lần 1.
\begingroup\renewcommand{\thetable}{\arabic{table}}% before chapter 1: no "0." prefix
\chapter*{Thông tin nhóm và phân công công việc}
\addcontentsline{toc}{chapter}{Thông tin nhóm và phân công công việc}

\section*{Thành viên}

Nhóm M9 gồm bảy học viên, đã chốt danh sách và đề tài 02-KH (phát hiện lỗi và định vị lỗi, hướng nghiên cứu khoa học).
Nhóm trưởng: \textbf{Doãn Văn An} (20252778M), đầu mối liên hệ với giảng viên và người giữ kho mã của nhóm.

\begin{table}[H]
\centering
\caption{Thành viên nhóm M9}
\label{tab:thanh-vien}
\small
\begin{tabularx}{\linewidth}{@{}r l l >{\raggedright\arraybackslash}X@{}}
\toprule
STT & Họ và tên & MSHV & Vai trò chính \\
\midrule
1 & Doãn Văn An & 20252778M & Nhóm trưởng; điều phối, hạ tầng thực nghiệm, phương pháp và viết bài \\
2 & Tạ Công Nam & 20252640M & Bộ nhãn và duyệt nhãn; phân tích kết quả \\
3 & Nguyễn Hoàng Hải & 20252583M & Tổng quan tài liệu và bộ từ khóa \\
4 & Đỗ Xuân Tráng & 20252774M & Lỗi thật BugsInPy (RQ3) \\
5 & Hoàng Trung Kiên & 20251333M & Biến thể đổi tên và các họ mô hình (RQ3) \\
6 & Mai Văn Tùng & 20252788M & Slide \\
7 & Nguyễn Đức Tùng & 20252772M & Công cụ phân tích tĩnh; kiểm tra số liệu \\
\bottomrule
\end{tabularx}
\end{table}

\section*{Phân công theo mảng việc}

Công việc của cả tiểu luận được chia thành bảy mảng lớn. Nhóm trưởng phụ trách phần lớn khối lượng (hạ tầng, tái hiện,
phương pháp đề xuất và viết bài); mỗi thành viên còn lại phụ trách chính một mảng và cùng tham gia duyệt nhãn.

\begin{table}[H]
\centering
\caption{Phân công công việc theo mảng}
\label{tab:phan-cong}
\small
\begin{tabularx}{\linewidth}{@{}>{\raggedright\arraybackslash}X >{\raggedright\arraybackslash}p{0.2\linewidth} >{\raggedright\arraybackslash}p{0.27\linewidth}@{}}
\toprule
Mảng việc & Phụ trách chính & Cùng thực hiện \\
\midrule
Tổng quan tài liệu, bộ từ khóa, phân tích ba bài báo được giao & Nguyễn Hoàng Hải & Doãn Văn An \\
Hạ tầng gọi LLM, tái hiện SemLoc và kiểm tra cách đánh giá (RQ1) & Doãn Văn An & Tạ Công Nam \\
Bộ nhãn: kiểm chứng bằng thực thi, rà soát nhãn, duyệt mù & Tạ Công Nam & Doãn Văn An; cả nhóm duyệt nhãn \\
Phương pháp đề xuất và so sánh (RQ2): prompt, thí nghiệm loại bỏ, các cách bổ sung, phân tích tĩnh & Doãn Văn An &
Nguyễn Đức Tùng \\
Độ bền (RQ3): lỗi thật BugsInPy, mã đổi tên, các họ mô hình & Đỗ Xuân Tráng, Hoàng Trung Kiên & Doãn Văn An \\
Phân tích thất bại, đóng góp, triết lý khoa học, mối đe dọa & Doãn Văn An & Tạ Công Nam, Hoàng Trung Kiên \\
Báo cáo, bài báo arXiv, slide, gói tái hiện & Doãn Văn An & Mai Văn Tùng, Nguyễn Đức Tùng \\
\bottomrule
\end{tabularx}
\end{table}

\section*{Công cụ AI sử dụng}

\begin{itemize}
\item \textbf{Claude Code} (Anthropic) là trợ lý chính: đọc và tóm tắt tài liệu, viết mã thực nghiệm và kiểm thử tự
động, chạy thực nghiệm, soạn báo cáo LaTeX và slide. Mọi thay đổi được nhóm trưởng duyệt trước khi đưa vào kho mã.
\item \textbf{LLM là đối tượng thực nghiệm}, gọi qua một proxy tương thích OpenAI: Gemini 3.8 Flash là mô hình chính;
Claude Sonnet 4.6 (mô hình của bài báo gốc) để đối chiếu khi tái hiện; cho RQ3 thêm Gemini 3.5 Flash-Lite,
Claude Sonnet 5.5, GPT-6 luna và gpt-oss-120b (mô hình mở của OpenAI).
\end{itemize}

Để kết quả kiểm chứng được dù có dùng AI, mọi con số trong báo cáo là macro sinh tự động từ dữ liệu, mọi lượt gọi LLM
được ghi nhật ký và lưu bộ đệm để chạy lại, và kho mã có bộ kiểm thử tự động.

\IfBelow{lan2}{%
\section*{Hiện trạng và đánh giá sau tuần làm việc đầu tiên}

Tính đến ngày 11/11/2026, một tuần sau khi chốt nhóm và đề tài, hiện trạng hai nội dung đầu theo hướng dẫn của môn học
như sau.

\begin{table}[H]
\centering
\caption{Hiện trạng sau tuần làm việc đầu tiên (11/11/2026)}
\label{tab:hien-trang}
\small
\begin{tabularx}{\linewidth}{@{}>{\raggedright\arraybackslash}p{0.22\linewidth} >{\raggedright\arraybackslash}X >{\raggedright\arraybackslash}p{0.17\linewidth}@{}}
\toprule
Nội dung & Hiện trạng & Tự đánh giá \\
\midrule
1. Tài liệu và từ khóa & Đã đọc ba bài báo được giao, tìm thêm tài liệu liên quan, lập bộ từ khóa và bảng so sánh
(Chương~\ref{chap:tong-quan}) & Đạt \\
2. RQ và thực nghiệm & Ba câu hỏi nghiên cứu, giả thuyết và thiết kế thực nghiệm (Chương~\ref{chap:phuong-phap}); hạ
tầng gọi LLM và tái hiện SemLoc; đã có kết quả RQ1 (Mục~\ref{sec:rq1}). RQ2 đang chạy, RQ3 chưa bắt đầu & Đang thực
hiện, đúng tiến độ \\
\bottomrule
\end{tabularx}
\end{table}

\textbf{Đánh giá chung.} Nhóm đã ổn định tổ chức, có hạ tầng thực nghiệm chạy lại được và kết quả đầu tiên: dòng hạng 1
của SemLoc luôn là dòng LLM tự đoán, và nhãn của bộ dữ liệu có vấn đề. \textbf{Khó khăn:} chạy lại SemLoc với chính mô
hình của bài báo không khớp nhật ký công bố; hạn mức gọi LLM của proxy giới hạn số lần chạy mỗi ngày; dựng môi trường
cho lỗi thật của BugsInPy tốn thời gian. \textbf{Hướng xử lý:} báo cáo cả hai cách (tính lại từ nhật ký và chạy lại),
lưu bộ đệm mọi lượt gọi LLM để không phải gọi lại, và dựng môi trường BugsInPy song song với RQ2.
}
\endgroup
